\documentclass[10pt]{article}

\usepackage[utf8]{inputenc}

\usepackage[T1]{fontenc}

\usepackage{amsmath}

\usepackage{amsfonts}

\usepackage{amssymb}

\usepackage[version=4]{mhchem}

\usepackage{stmaryrd}

\usepackage{bbold}



\begin{document}

ной функции Ляпунова. Частотные критерии У. а. обычно охватывают все критерии, к рые могут быть получены с помощью функций Ляпунова из нек-рых многопараметрич. классов функций, Напр., критериї (12) - необходимое и достаточное условие существования функции



\[

V(x)==x^{*} H x

\]



$\left(H=H^{*}=\mathrm{const}-(N \times N)\right.$-матрица, *- знак эрмитова сопряжения) такой, что ее производная в силу системы $(11)$, (2) с произвольной нелинейностью (2) (для к-рой $\left.\mu_{1} \leqslant \varphi(\sigma, t) / \sigma \leqslant \mu_{2}\right)$ удовлетворяет условию



\[

d V(x) / d t<0 \text { при } x \neq 0 .

\]



Аналогично, частотное условие Попова (13) охватывает все критерии, к-рые можно установить, используя функции Ляпунова вида



\[

V(x)=x^{*} I I x+\vartheta \int_{0}^{\sigma} \varphi(\sigma) d \sigma .

\]



Известно много других частотных критериев У. а. для разных классов нелинейностей (см. [3]-[6]). Они распространены, в частности, на важные для приложений случаи неединственного состояния равновесия (см. [1]). Частотные критерии У. а. позволили выделять классы нелинейных систем общего вида, для к-рых факт устойчивости в делом устанавливается особенно просто. Напр., для системы (1) с $\xi=\varphi(\sigma), \sigma=C x, N \leqslant 3, n=1$ (т. е. для произвольной системы не выше 3-го порядка с одной нелинейностью) имеет место асимптотическая устойчивость в целом, если $\mu_{1} \leqslant \varphi^{\prime}(\sigma) \leqslant \mu_{2}$ и если любая линейная система с $\xi=\mu \sigma, \mu_{1} \leqslant \mu \leqslant \mu_{2}$, асимптотически устойчива. Для систем 4-го (и более высокого) порядка аналогичное утверждение ошибочно. Более того, при $N \geqslant 4, n=1$ существуют такие системы (1) и нелинейности $\xi=\varphi(\sigma), \mu_{1} \leqslant \varphi^{\prime}(\sigma) \leqslant \mu_{2}$, что матрица любой линеаризованной систөмы с $\xi=\mu \sigma, \mu_{1} \leqslant \mu \leqslant \mu_{2},-$ матрица Гурвица, а нелинейная система имеет пөриодическое решение.



При замене условия минимальної устойчивости аналогичным условием минимальной неустойчивости неравенства (8), (11), (12), (13) становятся критериями а бс о Јі ю т о й н е у с т о й ч и в о с т И (при соответствующем понимании последнего термина). Пусть, напр., в действительном случае с $n=1$ матрица коэффициентов системы (1) с $\xi=\mu \sigma, \sigma=C x$ (т. е. матрица $A+$ $+B \mu C)$ при нск-ром $\mu, \mu_{1} \leqslant \mu \leqslant \mu_{2}$, имеет $k \geqslant 1$ собственных значений в полуплоскости $\operatorname{Re} \lambda>0$ и выполнено частотное условие (12). Тогда у системы (1), (2) с функцией $\varphi(\sigma, t)$, удовлетворяющей условию $\mu_{1} \leqslant \varphi(\sigma, t) / \sigma \leqslant \mu_{2}$ (а также у системы (1), (3)), имеются решения $x(t)$, для к-рых



\[

|x(t)| \geqslant \mathrm{Ce}^{\varepsilon t}|x(0)| \text { при } t \geqslant 0,

\]



где постоянные $C>0, \varepsilon>0$ - одни и те же для всех систем указанного класса. Соответствующие векторы $x(0)$ заполняют конус $x(0)^{*} H x(0)<0$, где $H=H^{*}-$ нек-рая матрица, имеющая $k$ отрицательных собственных значений.



Аналогично, пгр $n=1$ условие (13) является ч а стотны м к р итериемабсолют о ой не у ст о й ч и в о с т и системы (1) с $\sigma=C x$ в классе стационарных нелинейностей $\xi(t)=\varphi[\sigma(t)]$, где $0 \leqslant \sigma \varphi(\sigma) \leqslant$ $\leqslant \mu_{0} \sigma^{2}$, если в (1) матрица $A$ имеет собственные значения в полуілоскости $\operatorname{Re} \lambda>0$.



В теории У. а. установлены также аналогичные частотные критерии диссипативности, конвергенции, существования периодич. движений (автоколебаний и вынужденных режимов) и др. (см., нашр., [3], [5] и литературу в $[1],[3],[5])$.



Jum.: [1] Ге лиг А. X., ЈІ е о нов Г. А., Як у бов и ч В. А., Устойчивость нелинейных систем с неединственным состоянием равновесия, М., 1978; [2] А й з е р м а н М. А., Г а н т м а х е р Ф. Р., Абсолютная устойчивость нелинейных

регулируемых систем, М., $1963 ;[3]$ я к у б о в и ч В. А., в кн.: регулируемых систем, М., 1963; [3] Я к у б о в и ч В. А., в кн.: равления, М., 1975 , с. $74-180$; [4] П о п о в В. М., Гиперустойчивость автоматических систем, пср. с рум., М., 1970; [5] В оp о но в А. А., Устойчивость, управляемость, наблюдаемость, М., 1979; [6] Р е 3 в а н В., Аб̆солютная устойчивость автоматических систем с запаздыванием, пер. с рум., М., 1983; [7] П я т-

І и ц к и й Е. С., "Автоматика и телемеханика", 1968 , №, 6 , и и ц к и й Е. С., "Автоматика и 'телемеханика", 1968, Nㅡ 6,

с. $5-36$. УСТОЙЧИВОСТЬ ВЫЧИСЛИТЕЛЬНОГО АЛГОРИТМА - равномерная относительно $h$ и $m$ ограниченность частично разрешающих операторов $L_{m}^{h}$, описывающих последовательные аташы вычислительного алгоритма решения уравнения



\[

L^{h} u^{h}=f^{h},

\]



напр. сеточного уравнения с шагом $h$ (см. Замыкание вычислительного алгоритма). У. в. а. является гарантией слабого влияния вычислительной погрешности на результат вычислений. Однако не исключена возможность, что величина $p(h)=\sup _{m}\left\|L_{m}^{h}\right\|$ растет сравнительно медленно и соответствующее усиление влияния вычислительной погрешности при $h \rightarrow 0$ оказывается практически допустимым. Понятие У. в. а. конкретизируется в применении к проекционно-сеточным методам (см. [4]) и в применении к итерационным методам (см. [6]). Имеются и другие определения У. в. а. (см., напр., [1], [3]).



Лит.: [1] Б а б у шк а И., В и т а с с к Э., П р а ге р М., Численные процессы решения дифференциальных уравнений, пер. с англ., М., 1969; [2] Б а х в а л о в Н. С., Численные методы, 2 изд., М., 1975; [3] Га в у р и н М. К., Лекций по метоцам вычислений, М., 1971; [4] М а р ч у к Г. И., А г о ш[5] С а м а р с к и й А. А., Г у л и н А. В., Устойчивость разностных схем, М., 1973; [6] С а м а р с к и й А. А., Н и кол а с в Е. С., Методы решения сеточных уравнений, М., 1978.



УСТОЙЧИВОСТЬ ВЫЧИСЛИТЕЛЬНОГО ПРОЦЕССA - свойство, характеризующее скорость накопления суммарной вычислительной погрешности. Понятие У. в. п. было введено потому, что в реальных расчетах невозможно оперировать с точными числами, ибо нельзя избежать округления, к-рое иногда может быть причиной быстрой потери точности.



Вычислительный процесс есть последовательность арифметич. операций над числами. Пусть $X_{i}$-нормированное линейное пространство и $A_{i}$-непрерывный оператор $A_{i}: X_{1} \times X_{2} \times \ldots \times X_{i} \rightarrow X_{i+1}$. Тогда последовательность уравнений



\[

\begin{gathered}

x_{i+1}=A_{i}\left(x_{1}, x_{2}, \ldots, x_{i}\right), \\

x_{i+1} \in X_{i+1}, i=1,2, \ldots, N-1,

\end{gathered}

\]



задает вычислительный процесс с исходным данным $x_{1}$ и п ромежуточными результатами $x_{i}, i=2,3, \ldots, N-1$. Обытно $X_{i}=\mathbb{R}^{n}$, а оператор $A_{i}$ состоит из конечного числа арифметич. операций. Как правило, $x_{i+1}$ зависит не от всех ранее полученных промежутотных результатов. Число $N$ может быть задано заранее или определено в самом вычислительном процессе. В последнем слутае $N$ зависит от $x_{1}$ (напр., если $N$-число итераций, необходимых для достижения заданной точности).



Реальный вытислитөльный процесс не может быть п роведен в точном соответствии с определением (1), так как при выполнении ариф্фметич. операций допускаются ошибки округления и $x_{i+1}$ полутается из неточных предыдущих результатов. Это значит, что вместо әлемента $x_{i+1}$ фактически вычисляется әлемент



\[

\tilde{x}_{i+1}=A_{i}\left(\tilde{x}_{1}, \tilde{x}_{\mathbf{2}}, \ldots, \tilde{x}_{i}\right)+\delta_{i}, \delta_{i} \in X_{i+1} \text {, }

\]



где малая аддитпвная ошибка $\delta_{i}$ возникает из-за округлений в ходе выполнения оператора $A_{i}$. Значение $\delta_{i}$ оп ределяется значениями $\tilde{x}_{j}, j=1,2, \ldots, i$, спосо-





\end{document}